\section{Introduction}
\label{sec:intro}

The rapid evolution of modern manufacturing toward Industry 4.0 has significantly increased the complexity and integration of industrial systems. Within these sophisticated environments, the reliability and safety of critical equipment, such as rotating machinery and chemical reactors, are paramount to ensuring continuous production and preventing catastrophic failures. Industrial fault diagnosis (IFD) has thus emerged as a vital research area, aiming to detect and identify equipment anomalies through the analysis of sensor data, such as vibration signals, pressure, and temperature. Effective IFD systems not only reduce maintenance costs by enabling predictive maintenance strategies but also enhance human safety and environmental protection by preventing accidents caused by equipment breakdown. Recent years have seen the widespread adoption of deep learning-based methods for fault diagnosis, which have demonstrated remarkable performance in extracting discriminative features from high-dimensional industrial data. However, the efficacy of these supervised learning models heavily depends on the availability of massive labeled datasets covering all possible operating conditions and fault types, a requirement that is increasingly difficult to satisfy in real-world scenarios.

Despite the success of deep learning in IFD, the practical application of these models is often hampered by the severe challenge of data scarcity. In real-world industrial settings, collecting a comprehensive dataset that includes all potential fault types is inherently impractical, if not impossible. While data for normal operating conditions and common faults are relatively abundant, samples for rare or novel failure modes are extremely scarce. This imbalance is particularly pronounced for critical infrastructure where faults occur infrequently, or for newly deployed machinery that has not yet experienced its full range of degradation patterns. Furthermore, the cost of labeling industrial data is high, requiring domain expertise and often involving destructive testing to obtain specific fault samples. Consequently, supervised models trained on a limited set of "seen" classes often fail to generalize to "unseen" fault categories encountered during deployment. This limitation underscores the urgent need for a more robust diagnostic framework capable of identifying novel faults for which no training samples were previously available, leading to the rise of zero-shot learning (ZSL) in the industrial domain.

To address the absence of training data for unseen classes, zero-shot learning (ZSL) has been introduced into fault diagnosis by leveraging auxiliary information to bridge the gap between seen and unseen domains. Early ZSL methods for IFD primarily relied on attribute-based approaches, where each fault category is represented by a high-dimensional attribute vector. These attributes, often defined by domain experts, describe the physical characteristics of the fault, such as "increased vibration at high frequency" or "irregular oil pressure." By learning a compatibility function or a mapping between the feature space (extracted from sensor signals) and the attribute space, the model can infer the labels of unseen classes based on their predefined attribute patterns. Traditional frameworks like Direct Attribute Prediction (DAP) and Structured Joint Embedding (SJE) paved the way for this paradigm. However, these methods typically utilize binary attribute vectors (0 or 1), which provide a limited and rigid representation of the complex physical phenomena associated with industrial faults, often failing to capture the subtle nuances and continuous variations inherent in different failure stages.

The limitation of static embeddings led researchers to explore generative ZSL approaches, which utilize Generative Adversarial Networks (GANs) or more recently, Diffusion Models, to synthesize data for unseen classes. By training on seen classes, these models learn to generate synthetic features conditioned on class-level attributes. Once trained, the generator can produce artificial samples for unseen classes by taking their corresponding attribute vectors as input. These synthetic samples are then used to train a standard supervised classifier, effectively converting the ZSL problem into a conventional classification task. In the context of industrial systems, GAN-based models like f-CLSWGAN have been applied to augment fault data. More recently, Denoising Diffusion Probabilistic Models (DDPMs) have shown superior stability and diversity in generating high-fidelity industrial features. These generative methods significantly alleviate the seen-unseen domain shift by providing the classifier with a direct look at the distribution of unseen classes, yet they still heavily rely on the quality and semantic depth of the input conditioners (attributes).

Despite these advancements, existing ZSL methods for fault diagnosis suffer from several fundamental limitations that hinder their performance in complex industrial environments. First, the reliance on human-defined binary attributes is inherently flawed; such attributes lack semantic richness and are prone to subjectivity and errors during the labeling process. This "semantic gap" makes it difficult for the model to learn a robust mapping that generalizes to truly novel faults. Second, the domain shift problem—where the distribution of generated features for unseen classes differs from the real, latent distribution—remains a major obstacle. Most generative models treat the semantic space as a discrete set of points, ignoring the continuous manifold structure of fault characteristics. Consequently, the generated features often lack the fine-grained discriminative power required to distinguish between similar unseen faults. Lastly, current diffusion-based augmentation techniques often fail to maintain semantic-feature alignment during the denoising process, leading to a degradation in the global generalization capability of the diagnostic system across different industrial datasets.

In this paper, we propose a novel framework termed CLIP-Enhanced Semantic Manifold Diffusion (CESM-Diff) to overcome these challenges. Our approach moves beyond rigid binary attributes by integrating the Contrastive Language-Image Pretraining (CLIP) text encoder to extract continuous, high-dimensional semantic embeddings from natural language fault descriptions. This allows us to capture the rich, nuanced relationships between different fault types using the powerful linguistic knowledge base embedded in CLIP. To leverage these continuous embeddings, we design a dual-path diffusion architecture consisting of a Feature Path for denoising fault representations and a Semantic Path that simultaneously refines the semantic context. A cross-path attention mechanism is introduced to ensure tight alignment between the sensor-derived features and the textual semantics throughout the generative process. Furthermore, we introduce Semantic Manifold Interpolation (SMI), a technique that performs smooth interpolation within the CLIP semantic space to cover the latent manifold of unseen faults, thereby generating more realistic and discriminative features for novel failure modes.

The primary contributions of this work are summarized as follows:
\begin{enumerate}
    \item We are the first to integrate the CLIP text encoder into the zero-shot fault diagnosis pipeline, enabling the use of continuous and flexible natural language descriptions instead of restrictive binary attribute vectors for semantic modeling.
    \item We propose a novel dual-path semantic diffusion architecture that incorporates a bidirectional cross-path attention mechanism, ensuring superior alignment between the generated fault features and their underlying semantic meanings.
    \item We introduce the Semantic Manifold Interpolation (SMI) method, which leverages the continuous nature of the CLIP latent space to synthesize high-fidelity features for unseen faults, effectively modeling the transition between known and unknown failure states.
    \item Robustness and generalization are validated through comprehensive evaluations on three diverse industrial datasets: the Tennessee Eastman Process (TEP), a Hydraulic System monitoring dataset, and the CWRU Bearing dataset, where CESM-Diff consistently outperforms state-of-the-art ZSL and generative methods.
\end{enumerate}

The remainder of this paper is organized as follows. Section~\ref{sec:related} reviews the related literature on zero-shot learning, diffusion models, and CLIP applications in industrial contexts. Section~\ref{sec:method} provides a detailed description of the proposed CESM-Diff framework, including the CLIP-based semantic extraction and the dual-path diffusion process. Section~\ref{sec:experiments} presents the experimental setup, results, and comparative analysis across multiple benchmarks. Finally, Section~\ref{sec:conclusion} concludes the paper and discusses future research directions.
